\section*{Author preliminaries}
We denote $\mathbf x=(x_1,\ldots,x_n)$. $[n]$ denotes $\{1,\ldots,n\}$. For a polynomial $f\in\mathbb F[\mathbf x]$, we denote up to degree-$d$ part as $\mathrm{Hom}_{\le d}f$ and $|f|_0$ as the sparsity or the number of monomials in $f$. For a differentiable function $f(x)$, we denote $f^{(k)}(x):=d^kf/dx^k$, as the $k$-th derivative of $f$. We also recall the definition of gcd of two polynomials $f,g$ in the ring $\mathbb F[\mathbf x]$: $\gcd(f,g)=:h\iff h\mid f,h\mid g,\text{ and }h'\mid f,g\implies h'\mid h$. It is unique up to constant multiples.
\paragraph{Field.} We denote the underlying field as $\mathbb F$ and assume that it is algebraically closed. All our results hold when the characteristic is large or not algebraically closed, as we can go to polynomial extensions and work with it.
\paragraph{div and mod operations.} For polynomials $f$ and $g\in\mathbb F[\mathbf x]$, if $f=g\cdot h+r$, where $h,r\in\mathbb F[\mathbf x]$ such that $\deg(r)<\deg(g)$, then $h$ is called the quotient, denoted $f\mathbin{\mathrm{div}}g$, and $r$ is called the remainder, denoted $f\mathbin{\mathrm{mod}}g$. Operation $\mathrm{mod}$ may not be well-defined in the multivariate settings, however, if one assumes $g$ to be monic in a variable say $x_n$, it is always well-defined (by thinking $g$ to be a univariate in $x_n$). A polynomial $g$ is monic in $x_n$ if the leading coefficient (the nonzero coefficient of highest degree) of $x_n$ is a non-zero constant in $\mathbb F$. Of course, if $g\mid f$, then $f\mathbin{\mathrm{div}}g=h$ and $f\mathbin{\mathrm{mod}}g=0$, irrespective of monic-ness.
\paragraph{Circuit complexity.} For a polynomial $f\in\mathbb F[\mathbf x]$, we define the complexity $L_{\mathbb F}(f)$ of $f$ to be the length of the smallest division-free arithmetic circuit (with only $\{+,-,\times\}$ gates) computing $f$. We also define the complexity $\tau_{\mathbb F}(f)$ of $f$ to be the length of the smallest division and constant-free arithmetic circuit computing $f$ (all the constants are made from 1), for formal definition see Section 5.2 of the original paper. We will remove subscript $\mathbb F$ when the underlying field is clear from the context.
\paragraph{Border complexity context (author Section 3.3).} The notion of border (equivalently, approximative) complexity is important in computer science. This concept popped up from early works on matrix multiplication and border rank of tensors, see [10]). Whether approximation of polynomials provides any additional computational power is a natural question which fundamentally motivated the foundation of Geometric Complexity theory (GCT). The notion of border complexity can be motivated through two ways: topological and algebraic, and both the perspectives are known to be equivalent [1]. For further details, we refer to [20,36]. In this paper, we only work with algebraic approximation upper bounds. In the algebraic definition, one can talk about the convergence $\epsilon\to0$. Here, one can see $\epsilon$ as a formal variable and $\mathbb F(\epsilon)$ as the function field. For an algebraic complexity class $\mathcal C$, the approximation is defined as follows [7,Definition2.1].
\paragraph{Definition 20 (Approximative closure of a class [7]).} Let $\mathcal C$ be an algebraic complexity class over field $\mathbb F$. A family $(f_n)$ of polynomials from $\mathbb F[\mathbf x]$ is in the class $\overline{\mathcal C}(\mathbb F)$ if there are polynomials $f_{n;i}$ and a function $t:\mathbb N\mapsto\mathbb N$ such that $g_n$ is in the class $\mathcal C$ over the field $\mathbb F(\epsilon)$ with $g_n(\mathbf x)=f_n(\mathbf x)+\epsilon f_{n,1}(\mathbf x)+\epsilon^2f_{n,2}(\mathbf x)+\cdots+\epsilon^{t(n)}f_{n,t(n)}(\mathbf x)$.
\paragraph{Definition 21 ([8,Defn.3.1]).} Let $f\in\mathbb F[\mathbf x]$. The border complexity $\underline L(f)$ is the smallest number $r$, such that there exists $F$ in $\mathbb F(\epsilon)[\mathbf x]$ satisfying $F|_{\epsilon=0}=f$ and $L_{\mathbb F(\epsilon)}(F)\le r$.
Note that the circuit of $F$ may be using $1/\epsilon$ in an intermediate step. So, we cannot merely assign $\epsilon=0$ and get a $\epsilon$-free circuit. Also, the $\epsilon$-degree can be exponential in its size (and thus cannot be interpolated), see [8,Theorem5.7]). Thus, potentially $\underline L(f)$ can be significantly smaller than $L(f)$. The above definition can be used to define closures of complexity class, e.g., $\overline{\mathrm{VP}}$. In this case, one can assume wlog that the degrees of $g_n$ and $f_{n,i}$ are $\mathrm{poly}(n)$. It is known to be closed under factoring [8,Theorem4.1]. However, the usual method of Hensel-lifting does not work when the given circuit class computes polynomials of super-polynomial degree. Also, Strassen's method would have a dependency on the degree of the final polynomial. However, we can prove Theorem 18 analogously, in the border sense.
\paragraph{Lemma 61 (Coefficient-Extraction).} Coefficient-extraction in arithmetic circuits is easy using interpolation, see the folklore lemma below, for a proof see [43]. Let $L(f)=s$ with $f\in\mathbb F[\mathbf x]$ and $f=\sum_{0\le i\le d}f_ix_n^i$ with $f_i\in\mathbb F[x_1,\ldots,x_{n-1}]$. Then there is a circuit $C$ of size $O(sd^2)$ computing $f_0,f_1,\ldots,f_d$.
\paragraph{Lemma 62 (Homogenization).} The next lemma is a homogenization trick, used in [45]. For a proof, see [43,Theorem2.2]. If $f$ has an arithmetic circuit of size $s$, then for any $d$, there is a circuit of size $O(sd^2)$ computing $\mathrm{Hom}_{\le d}f$.
\paragraph{Lemma 64 (Polynomial Identity Lemma).} Let $p(\mathbf x)$ be an $n$-variate nonzero polynomial of degree $d$. Let $S\subseteq\mathbb F$ be a finite set. Then $\Pr_{\alpha\sim S^n}[p(\alpha)=0]\le d/|S|$. Here, $\alpha\in S^n$ is picked independently and uniformly at random. This is the standard Polynomial Identity Lemma [38,15,51,41].
